\startchapter[title={Detalles de las metaheurísticas}, reference=Anexo_Metaheuristicas]

Este anexo reúne detalles específicos de las dos metaheurísticas del
Capítulo~\in[Capitulo_4]. Estos detalles incluyen la construcción de las rutas, que comparten \ac{hc} y \ac{aco}
(Sección~\in[sec:anexo_rutas]), algunos aspectos de la implementación
(Sección~\in[sec:anexo_impl]), y las constantes de implementación y los
estudios con los que se evaluaron sus parámetros
(Secciones~\in[sec:calib_hc] y~\in[sec:calib_aco]). Los parámetros de
cada metaheurística, con sus valores por defecto, se presentan en las
Secciones~\in[sec:hc_parametros] y~\in[sec:aco_parametros].

%----------------------------------------------------------------------------------------
\startsection[title={Construcción de las rutas}, reference=sec:anexo_rutas]

Esta sección detalla las construcciones de la ruta inicial de cada lote en
\ac{hc} (Sección~\in[sec:hc_solucion_inicial]), con la notación de la
Sección~\in[sec:heur_rutas]. El pulido y la búsqueda de respaldo se usan también en \ac{aco}
(Secciones~\in[sec:aco_construccion] y~\in[sec:aco_pulido]).

%----------------------------------------------------------------------------------------
\startsubsection[title={Construcciones candidatas}, reference=sec:anexo_rutas_construcciones]

\startitemize[n]
	\item {\bf Vecino más cercano.} Desde el depósito, se va repetidamente a
	la ubicación pendiente más cercana por el camino más corto (algoritmo de
	Dijkstra) que no pasa por ubicaciones ya visitadas ni por
	$v_{\text{end}}$. Las ubicaciones requeridas por las que se pasa en el
	camino cuentan como visitadas. Cuando no quedan pendientes, se llega a
	$v_{\text{end}}$ por el camino más corto que evita lo ya recorrido. Es
	determinística, y falla si en algún paso ninguna ubicación pendiente es
	alcanzable.
	\item {\bf Inserción del más lejano.} Primero se decide el orden en que
	se visitan las ubicaciones requeridas con la inserción del más lejano
	\cite[rosenkrantz1977analysis]. Se parte de la secuencia
	$(v_{\text{start}}, v_{\text{end}})$ y, repetidamente, se toma la
	ubicación pendiente más alejada de las que ya están en la secuencia y se
	la inserta en la posición que menos alarga el recorrido. En este paso las
	distancias son las de los caminos más cortos en el grafo, sin considerar
	la restricción de no repetir ubicaciones. Después, ese orden se realiza
	uniendo cada par de ubicaciones consecutivas con el camino más corto que
	evita todo lo ya recorrido. Esta construcción se aplica dos veces: sobre
	el grafo original y sobre el grafo invertido, en el que se invierte el
	sentido de todos los arcos y se intercambian $v_{\text{start}}$ y
	$v_{\text{end}}$. Una ruta del grafo invertido, leída al revés, es una
	ruta del grafo original con la misma distancia, así que esto equivale a
	construir la ruta desde el final hacia el principio.
	\item {\bf Construcciones aleatorias.} Se generan 20 rutas con una regla
	aleatoria. Desde la ubicación actual, cada ubicación pendiente $r$ que se
	puede alcanzar sin repetir ubicaciones recibe un peso $(1/d_r)^{\beta}$,
	donde $d_r$ es la distancia de ese camino y $\beta = 2{,}5$, y se elige
	una con probabilidad proporcional a su peso. Antes de aceptar la
	ubicación elegida, se verifica que la ruta todavía pueda completarse: cada
	ubicación pendiente debe estar en algún camino simple entre la ubicación
	elegida y $v_{\text{end}}$ que no pase por lo ya recorrido. Esto se
	comprueba de forma exacta con las componentes biconexas del grafo. La
	condición es necesaria pero no suficiente, así que la construcción
	todavía puede trabarse. Cuando ninguna candidata la cumple, se deshace el
	último paso y se prueba otra ubicación en su lugar, con un máximo de 30
	retrocesos. Las construcciones alternan entre el grafo original y el
	invertido, porque la regla favorece a las ubicaciones cercanas y deja para
	el final a las más alejadas del punto de partida. Cada ruta completa se
	mejora alternando dos búsquedas locales hasta que deja de acortarse:
	2-opt sobre las ubicaciones de la ruta (Ecuación~\in[eq:2opt]) y el
	re-ruteo de tramos, que reemplaza el tramo entre dos ubicaciones
	requeridas consecutivas por el camino más corto que evita el resto de la
	ruta, cuando este es más corto.
\stopitemize

Las rutas se comparan por distancia y no por tiempo de ruta $r_b$, que
también suma el tiempo de servicio de las ubicaciones de paso. Ambos
criterios coinciden cuando las ubicaciones no requeridas tienen $S_i = 0$,
como en todos los escenarios de estrés, y en las instancias funcionales en
que no es así la ruta más corta es también la de menor $r_b$.

Una única construcción no alcanza. Con solo el vecino más cercano mejorado con 2-opt, en el escenario de estrés
con grilla de 7×7 del perfil realista las rutas sumaban 490 unidades de
tiempo, contra 425 de las rutas óptimas de la Etapa~1 del modelo Split. En
el escenario con grilla de 14×14, un lote no conseguía ruta y la instancia
resultaba infactible, aunque el modelo Split la resolvía. Generar varias
soluciones de partida y mejorarlas con búsqueda local es una forma habitual
de diversificar el inicio de una búsqueda local; Coruzzolo et al. la usan
para el loteo \cite[coruzzolo2023orderpicking].

\stopsubsection

%----------------------------------------------------------------------------------------
\startsubsection[title={Pulido}, reference=sec:anexo_rutas_pulido]

Una vez obtenidas las rutas candidatas, se selecciona la de menor distancia
y se aplica un pulido. El pulido repite los siguientes pasos hasta que la
ruta deja de acortarse, con un máximo de cinco rondas:
\startitemize
	\item {\bf 2-opt sobre las ubicaciones requeridas.} Se invierte el orden
	de un bloque de hasta 8 ubicaciones requeridas consecutivas y se vuelven a
	trazar los caminos entre ellas, evitando el resto de la ruta. A
	diferencia del 2-opt sobre todas las ubicaciones de la ruta, que en una
	grilla casi nunca es aplicable porque los dos nuevos arcos deben unir
	ubicaciones adyacentes, este movimiento redibuja los caminos completos.
	\item {\bf Reubicación de ubicaciones.} Es una versión simplificada del
	operador {\em or-opt}, propuesto por Or en su tesis doctoral de 1976
	\cite[babin2007oropt]. El operador original reubica cadenas de tres, dos
	y un vértice consecutivos; aquí se reubica una sola ubicación requerida
	por vez en otra posición de la secuencia, y se vuelven a trazar los
	caminos afectados.
	\item {\bf 2-opt y re-ruteo de tramos}, como en las construcciones
	aleatorias.
\stopitemize

\stopsubsection

%----------------------------------------------------------------------------------------
\startsubsection[title={Búsqueda de respaldo}, reference=sec:anexo_rutas_respaldo]

Si ninguna de las construcciones encuentra una ruta, se recurre a una
búsqueda exhaustiva de tipo {\em best-first}, sobre estados formados por la
ubicación actual, el conjunto de ubicaciones pendientes y el camino recorrido.
Los estados se exploran primero por cantidad de ubicaciones pendientes y
después por la distancia a la ubicación requerida más cercana, con un tope de
2.000.000 de expansiones. Si se alcanza el tope sin encontrar una ruta, el
lote queda sin ruta y la instancia se informa como infactible.

\stopsubsection

\stopsection

%----------------------------------------------------------------------------------------
\startsection[title={Otros detalles de implementación}, reference=sec:anexo_impl]

%----------------------------------------------------------------------------------------
\startsubsection[title={Peso de la penalización}, reference=sec:anexo_lambda]

El valor de $\lambda$ de la Ecuación~\in[eq:hc_costo] se eligió comparándolo
con cuánto puede aumentar $f$ cuando un lote tarda una unidad de tiempo más.
En el peor caso, ese lote y todos los que le siguen en la secuencia de su
{\em picker} terminan una unidad más tarde, y el adelanto o la demora de cada
una de sus órdenes cambia en una unidad, de modo que el aumento de $f$ está
acotado por
\startformula
	K = W^{c} \, |\mathcal{B}| + W^{\text{pen}} \max\left(W^{e}, W^{t}\right) |\mathcal{N}|.
\stopformula
Con $\lambda > K$, una unidad de exceso sobre $T^{\text{cold}}$ pesa en $f'$
más de lo que una unidad de tiempo adicional en un lote puede sumar a $f$,
de modo que reducir el exceso tiene prioridad sobre mejorar $f$. A
diferencia de las constantes M grande de los modelos lineales, un valor alto
de $\lambda$ no trae problemas numéricos, porque se usa solo para comparar
dos costos y no dentro de un resolvedor.

\stopsubsection

%----------------------------------------------------------------------------------------
\startsubsection[title={Evaluación incremental en \ac{hc}}, reference=sec:anexo_incremental]

Cada vecino difiere de la solución actual en pocos lotes, así que $f$ y $v$
no se recalculan desde cero, sino a partir de los valores de la última
solución aceptada. Para cada {\em picker}, los instantes de finalización se
recalculan desde la primera posición de su secuencia en la que cambió el lote
o su ruta, ya que las posiciones anteriores no se ven afectadas
(Ecuación~\in[eq:hc_finalizacion]). El adelanto y la demora se recalculan
solo para las órdenes de los lotes cuyo instante de finalización cambió, y el
exceso sobre $T^{\text{cold}}$, solo para los lotes cuya ruta o cuyo
{\em picker} cambió, porque no depende de la posición en la secuencia. Los
valores guardados se actualizan únicamente cuando el vecino es aceptado.

\stopsubsection

%----------------------------------------------------------------------------------------
\startsubsection[title={Control del tiempo}, reference=sec:anexo_tiempo]

En \ac{hc}, el tiempo se mide desde el comienzo de la ejecución, incluyendo
la verificación previa y la construcción de la solución inicial, y se
controla al inicio de cada iteración. Por lo tanto, no interrumpe la
construcción ni una iteración en curso.

En \ac{aco}, el tiempo se controla antes de construir cada hormiga, así que
una iteración puede cortarse a mitad de camino. Las hormigas dejan de
construirse al consumir la mitad de $\text{Max}_{\text{time}}$, pero solo si
ya existe una solución factible, de esa iteración o de una anterior, para
que el corte no produzca un falso infactible. La otra mitad queda para el
pulido y la búsqueda local sobre la asignación, que respetan el plazo
completo. Por qué el tiempo se reparte de esta forma se explica en la
Sección~\in[sec:calib_aco_tiempo].

El pulido también respeta el tiempo máximo. Si en una iteración hay varios
récords para pulir, cada uno recibe una parte igual del tiempo que queda,
reservando una parte más para la búsqueda local sobre la asignación. Al
vencer su plazo, el pulido se detiene y devuelve la mejor ruta obtenida
hasta ese momento, que nunca es más larga que la original.

\stopsubsection

%----------------------------------------------------------------------------------------
\startsubsection[title={Intentos y ruta de respaldo en \ac{aco}}, reference=sec:anexo_aco_respaldo]

Cada hormiga hace hasta 10 intentos de construcción por lote, alternando la
dirección entre intentos y empezando por una elegida al azar. Si los 10
fallan, se usa una ruta de respaldo determinística. Primero se prueba la
inserción del más lejano sobre el grafo original y, si no se puede realizar,
se prueba sobre el invertido. Si ninguna funciona, se usa la búsqueda
{\em best-first} con el tope de 2.000.000 de expansiones, como en \ac{hc}
(Sección~\in[sec:anexo_rutas_respaldo]). La ruta de respaldo se pule y se
guarda, de modo que se calcula una sola vez por lote y no una vez por
hormiga. Si tampoco se encuentra una ruta de respaldo, la hormiga se
descarta.

\stopsubsection

\stopsection

%----------------------------------------------------------------------------------------
\startsection[title={Calibración de Hill Climbing}, reference=sec:calib_hc]

%----------------------------------------------------------------------------------------
\startsubsection[title={Constantes de la implementación}, reference=sec:calib_hc_constantes]

El Cuadro~\in[tab:hc_constantes] reúne las constantes de \ac{hc}, que no se
configuran en cada ejecución. Se fijaron a partir de las mediciones
descritas en el Capítulo~\in[Capitulo_4], y son las mismas en los escenarios
funcionales y en los de estrés.

\startplacetable[reference=tab:hc_constantes, location=here, title={Constantes de implementación de \ac{hc}.}]
\bTABLE[split=yes, align=middle, frame=on, width=fit]
\setupTABLE[row][1][background=color, backgroundcolor=darkgray, style=bold]
\setupTABLE[column][3][width=8cm, align=flushleft]
\bTR \bTH Constante \eTH \bTH Valor \eTH \bTH Función \eTH \eTR
\bTR \bTD $\lambda$ \eTD \bTD 1000 \eTD \bTD Peso de la penalización por cadena de frío (Ecuación~\in[eq:hc_costo]). \eTD \eTR
\bTR \bTD Construcciones aleatorias \eTD \bTD 20 \eTD \bTD Rutas aleatorias generadas por lote en la solución inicial. \eTD \eTR
\bTR \bTD $\beta$ \eTD \bTD 2,5 \eTD \bTD Exponente de la distancia en la regla de elección de las construcciones aleatorias. \eTD \eTR
\bTR \bTD Retrocesos \eTD \bTD 30 \eTD \bTD Pasos que puede deshacer cada construcción aleatoria antes de abandonarse. \eTD \eTR
\bTR \bTD Bloque de 2-opt \eTD \bTD 8 \eTD \bTD Ubicaciones requeridas por bloque invertido en el 2-opt del pulido. \eTD \eTR
\bTR \bTD Rondas de pulido \eTD \bTD 5 \eTD \bTD Rondas máximas del pulido de la ruta inicial. \eTD \eTR
\bTR \bTD Expansiones {\em best-first} \eTD \bTD 2.000.000 \eTD \bTD Tope de la búsqueda exhaustiva de respaldo. \eTD \eTR
\bTR \bTD $\varepsilon$ \eTD \bTD $10^{-9}$ \eTD \bTD Tolerancia de la cadena de frío, la misma de los modelos lineales. \eTD \eTR
\eTABLE
\stopplacetable

\stopsubsection

%----------------------------------------------------------------------------------------
\startsubsection[title={Presupuesto de búsqueda}, reference=sec:calib_hc_presupuesto]

Para evaluar si los valores de estrés de $\text{Max}_{\text{it}}$ y
$\text{Max}_{\text{imp}}$ (Cuadro~\in[tab:hc_parametros]) limitan la
búsqueda, se repitieron los escenarios de estrés 02 a 19 con 200 y 120, 500
y 60, y 500 y 120 iteraciones máximas e iteraciones sin mejora,
respectivamente, con cinco semillas por escenario. El 01 queda afuera porque
la búsqueda no llega a iterar (Sección~\in[sec:heur_limitaciones]). Con 200 y
60, la búsqueda se detiene por estancamiento en 82 de las 88 ejecuciones que
encuentran una solución factible, y solo alcanza el tope de 200 iteraciones
en el 03, con todas las semillas, y en el 18, con una. Los presupuestos
mayores nunca empeoran la solución y la mejoran solo en cuatro escenarios:
el 03, con un único {\em picker} y 40 lotes para secuenciar, mejora hasta un
0,59~\%; el 15, el 17 y el 18, como máximo un 0,11~\%. Las diferencias de
tiempo de ejecución son del orden del ruido de la medición, de hasta 2~s en
uno u otro sentido, porque el tiempo lo domina la construcción de la
solución inicial y cada iteración es barata. Por lo tanto, con 200 y 60 la
búsqueda termina, en casi todos los escenarios, porque deja de encontrar
mejoras y no por el tope de iteraciones.

\stopsubsection

%----------------------------------------------------------------------------------------
\startsubsection[title={Variación entre semillas}, reference=sec:calib_hc_semillas]

Al cambiar los sorteos en cinco ejecuciones de cada escenario de estrés, con
200 iteraciones máximas y 60 sin mejora, la diferencia entre el mejor y el
peor valor objetivo fue menor que un 0,1~\% en todos los escenarios salvo el
17 (0,12~\%) y el 03 (0,44~\%), con una mediana de 0,02~\%. Además, con una
de las cinco semillas, dos escenarios terminaron sin solución factible: en
el 11, que tiene un límite de cadena de frío efectivo, la búsqueda no
alcanzó una solución que lo cumpla, y en el 09 no se consiguió construir la
ruta inicial de un lote. En ambos casos el resultado fue el mismo con los
cuatro presupuestos de búsqueda (Sección~\in[sec:heur_limitaciones]).

\stopsubsection

\stopsection

%----------------------------------------------------------------------------------------
\startsection[title={Calibración de Ant Colony Optimization}, reference=sec:calib_aco]

%----------------------------------------------------------------------------------------
\startsubsection[title={Constantes de la implementación}, reference=sec:calib_aco_constantes]

El Cuadro~\in[tab:aco_constantes] reúne las constantes de \ac{aco}, que no
se configuran en cada ejecución y son las mismas en los escenarios
funcionales y en los de estrés.

\startplacetable[reference=tab:aco_constantes, location=here, title={Constantes de implementación de \ac{aco}.}]
\bTABLE[split=yes, align=middle, frame=on, width=fit]
\setupTABLE[row][1][background=color, backgroundcolor=darkgray, style=bold]
\setupTABLE[column][3][width=8cm, align=flushleft]
\bTR \bTH Constante \eTH \bTH Valor \eTH \bTH Función \eTH \eTR
\bTR \bTD $\lambda$ \eTD \bTD 1000 \eTD \bTD Peso de la penalización por cadena de frío (Ecuación~\in[eq:hc_costo]). \eTD \eTR
\bTR \bTD $\tau_{\max} / \tau_{\min}$ \eTD \bTD 50 \eTD \bTD Relación entre las cotas de la feromona de ruta (Ecuación~\in[eq:aco_cotas]). \eTD \eTR
\bTR \bTD $\tau^{P}$ inicial / mínima \eTD \bTD 1 / 0,01 \eTD \bTD Valor inicial y mínimo de la feromona de asignación. \eTD \eTR
\bTR \bTD Peso elitista \eTD \bTD 2 \eTD \bTD Multiplicador del depósito de $s^*$ en la feromona de asignación. \eTD \eTR
\bTR \bTD Intentos por lote \eTD \bTD 10 \eTD \bTD Construcciones por lote y por hormiga antes de usar la ruta de respaldo. \eTD \eTR
\bTR \bTD Retrocesos \eTD \bTD 30 \eTD \bTD Pasos que puede deshacer cada construcción antes de abandonarse. \eTD \eTR
\bTR \bTD Bloque de 2-opt \eTD \bTD 8 \eTD \bTD Ubicaciones requeridas por bloque invertido en el pulido. \eTD \eTR
\bTR \bTD Rondas de pulido \eTD \bTD 5 \eTD \bTD Rondas máximas del pulido de cada récord. \eTD \eTR
\bTR \bTD Expansiones {\em best-first} \eTD \bTD 2.000.000 \eTD \bTD Tope de la búsqueda exhaustiva de respaldo. \eTD \eTR
\bTR \bTD Fracción para hormigas \eTD \bTD 0,5 \eTD \bTD Parte de $\text{Max}_{\text{time}}$ que pueden usar las hormigas; el resto queda para el pulido y la búsqueda local. \eTD \eTR
\bTR \bTD $\varepsilon$ \eTD \bTD $10^{-9}$ \eTD \bTD Tolerancia de la cadena de frío, la misma de los modelos lineales. \eTD \eTR
\eTABLE
\stopplacetable

\stopsubsection

%----------------------------------------------------------------------------------------
\startsubsection[title={Reparto del tiempo máximo}, reference=sec:calib_aco_tiempo]

En la grilla de 100×100 con $\text{Max}_{\text{time}} = 600$~s, dejar que
las hormigas usaran todo el tiempo dejaba 11~s de pulido y nada para la
búsqueda local sobre la asignación, con un objetivo de 1.403.120, contra
917.328 cuando el posprocesamiento se completaba. Por eso las hormigas
dejan de construirse al consumir la mitad de $\text{Max}_{\text{time}}$
(Sección~\in[sec:aco_parametros]).

\stopsubsection

%----------------------------------------------------------------------------------------
\startsubsection[title={Contribución de la búsqueda local}, reference=sec:calib_aco_ablacion]

Para medir cuánto del resultado aporta la colonia y cuánto las búsquedas
locales, se ejecutaron tres variantes en los escenarios de estrés 04 (grilla
de 7×7), 06 (grilla de 14×14) y 15 (6 lotes y 9 {\em pickers}), con tres
semillas cada uno: solo la colonia, sin ninguna búsqueda local; la colonia
con la búsqueda local sobre cada ruta, el pulido y la recombinación; y la
versión completa, que agrega la búsqueda local sobre la asignación. El
Cuadro~\in[tab:aco_ablacion] muestra la brecha respecto del objetivo del
modelo Split.

\startplacetable[reference=tab:aco_ablacion, location=here, title={Brecha de \ac{aco} respecto del modelo Split, en porcentaje, por variante (mínimo y máximo sobre tres semillas).}]
\bTABLE[split=yes, align=middle, frame=on, width=fit]
\setupTABLE[row][1][background=color, backgroundcolor=darkgray, style=bold]
\bTR \bTH Escenario \eTH \bTH Solo colonia \eTH \bTH + rutas \eTH \bTH Completa \eTH \eTR
\bTR \bTD 04 (7×7) \eTD \bTD 0,008 -- 0,013 \eTD \bTD 0,007 -- 0,012 \eTD \bTD 0,000 -- 0,002 \eTD \eTR
\bTR \bTD 06 (14×14) \eTD \bTD 0,025 -- 0,040 \eTD \bTD 0,025 -- 0,047 \eTD \bTD 0,010 -- 0,030 \eTD \eTR
\bTR \bTD 15 \eTD \bTD 0,026 -- 0,046 \eTD \bTD 0,008 -- 0,019 \eTD \bTD 0,005 -- 0,008 \eTD \eTR
\eTABLE
\stopplacetable

La colonia sola ya queda a menos de un 0,05~\% del modelo Split. Las
búsquedas locales sobre las rutas reducen la brecha sobre todo en el
escenario 15, donde la suma de los tiempos de ruta baja de entre 685 y 712 a
entre 672 y 676; en el 04 y el 06 su efecto queda dentro de la variación
entre semillas. La búsqueda local sobre la asignación reduce la brecha en
los tres escenarios, y es la que más la reduce en el 04 y el 06. Con la
versión completa, en el escenario 04 una de las semillas alcanza el objetivo
del modelo Split. Los tiempos de ejecución de las tres variantes son
similares, porque los domina la construcción de las hormigas: unos 15~s en
el 04, entre 109 y 116~s en el 06 y entre 44 y 51~s en el 15.

\stopsubsection

%----------------------------------------------------------------------------------------
\startsubsection[title={Sensibilidad a $\alpha$, $\beta$ y $\rho$}, reference=sec:calib_aco_sensibilidad]

Para evaluar los valores de la colonia, se ejecutó el escenario 04 sin
búsquedas locales, con todas las combinaciones de
$\alpha \in \{0{,}5;\ 1;\ 2\}$, $\beta \in \{1;\ 2{,}5;\ 4\}$ y
$\rho \in \{0{,}05;\ 0{,}15;\ 0{,}30\}$, y tres semillas por combinación.
Los rangos abarcan los valores por defecto y los que reportan Alaya y Solnon
\cite[alaya2007antcolony] y Lesch et al. \cite[lesch2023optimizing]. La
brecha media fue de 0,011~\% con los valores por defecto, de 0,009~\% con la
mejor combinación ($\alpha = 2$, $\beta = 4$, $\rho = 0{,}05$) y de
0,082~\% con la peor ($\alpha = 0{,}5$, $\beta = 1$, $\rho = 0{,}05$). El
parámetro con más efecto es $\beta$: con $\beta = 1$, la brecha media fue de
al menos 0,018~\% en todas las combinaciones, mientras que con $\beta = 4$
no superó el 0,028~\%. Los valores por defecto quedan, por lo tanto, cerca
de los mejores, y las diferencias entre ellos son menores que las que aporta
la búsqueda local.

\stopsubsection

%----------------------------------------------------------------------------------------
\startsubsection[title={Variación entre semillas}, reference=sec:calib_aco_semillas]

Se ejecutaron los escenarios de estrés 02 a 19 con cinco semillas cada uno y
los valores por defecto (Cuadro~\in[tab:aco_parametros]). El 01 queda afuera
por su duración, de unos 32 minutos por ejecución, y para usar los mismos
escenarios que en \ac{hc} (Sección~\in[sec:calib_hc_semillas]). Las
ejecuciones se repartieron en tres procesos en paralelo, así que sus tiempos
no son comparables con los de la suite de estrés y no se informan.

Las 90 ejecuciones encontraron una solución factible. La diferencia entre el
mejor y el peor valor objetivo de cada escenario tuvo una mediana de
0,005~\%. En los escenarios 02, 04 y 10, las cinco semillas llegaron al mismo
objetivo, y el máximo fue de 0,053~\% en el 11, el único con la cadena de
frío activa; en el resto no superó el 0,016~\%. En los escenarios que el
modelo Split resuelve, la brecha respecto de su óptimo quedó entre 0 y
0,072~\%, también con el máximo en el 11.

La variación es, por lo tanto, un orden de magnitud menor que la de \ac{hc},
cuya mediana fue de 0,02~\% y su máximo de 0,44~\%, y a diferencia de
\ac{hc}, ninguna semilla dejó un escenario sin solución.

\stopsubsection

\stopsection
\stopchapter
